% Part: first-order-logic
% Chapter: sequent-calculus
% Section: structural-rules

\documentclass[../../../include/open-logic-section]{subfiles}

\begin{document}

\iftag{FOL}
      {\olfileid{fol}{seq}{srl}}
      {\olfileid{pl}{seq}{srl}}

\olsection{Regels die de volgorde en het aantal zinnen veranderen}

Er zijn ook regels nodig waarmee we !!{sentence}s links en rechts van
een sequent op een andere plek kunnen zetten. Een logische regel werkt
alleen als de !!{sentence}s waarop hij ingrijpt helemaal links of
helemaal rechts in zijn uitgangssequent staan. Daarom is er een regel
om !!{sentence}s te verwisselen. Met andere regels kunnen we twee
gelijke !!{sentence}s samenvoegen tot één of aan een van beide kanten
!!a{sentence} toevoegen.

\subsection{Een zin toevoegen (verzwakking)}

\begin{defish}
\Axiom$ \Gamma \fCenter \Delta $
\RightLabel{\LeftR{\Weakening}}
\UnaryInf$ !A, \Gamma \fCenter \Delta$
\DisplayProof
\hfill
\Axiom$ \Gamma \fCenter \Delta$
\RightLabel{\RightR{\Weakening}}
\UnaryInf$ \Gamma \fCenter \Delta, !A$
\DisplayProof
\end{defish}

\subsection{Twee gelijke zinnen samenvoegen (contractie)}

\begin{defish}
\Axiom$ !A, !A, \Gamma \fCenter \Delta $
\RightLabel{\LeftR{\Contraction}}
\UnaryInf$ !A, \Gamma \fCenter \Delta$
\DisplayProof
\hfill
\Axiom$ \Gamma \fCenter \Delta, !A, !A$
\RightLabel{\RightR{\Contraction}}
\UnaryInf$ \Gamma \fCenter \Delta, !A$
\DisplayProof
\end{defish}

\subsection{De volgorde veranderen (verwisseling)}

\begin{defish}
\Axiom$ \Gamma, !A, !B, \Pi \fCenter \Delta $
\RightLabel{\LeftR{\Exchange}}
\UnaryInf$ \Gamma, !B, !A, \Pi \fCenter \Delta$
\DisplayProof
\hfill
\Axiom$ \Gamma \fCenter \Delta, !A, !B, \Lambda$
\RightLabel{\RightR{\Exchange}}
\UnaryInf$ \Gamma \fCenter \Delta, !B, !A, \Lambda$
\DisplayProof
\end{defish}

Een reeks stappen met verzwakking, contractie en verwisseling wordt
vaak weergegeven met dubbele afleidingslijnen.

De volgende regel heet de ``snederegel''. Strikt genomen kunnen we
zonder deze regel, maar hij maakt het veel makkelijker om bestaande
!!{derivation}s opnieuw te gebruiken en met elkaar te combineren.
\begin{defish}
\[
\Axiom$ \Gamma \fCenter \Delta, !A$
\Axiom$ !A, \Pi \fCenter \Lambda $
\RightLabel{\Cut}
\BinaryInf$ \Gamma, \Pi \fCenter \Delta, \Lambda$
\DisplayProof
\]
\end{defish}
\end{document}
